\section{Introduction}

We prove worst-start total-variation cutoff for single-site heat-bath
dynamics throughout the high-temperature phase of the zero-field
Gaussian Sherrington--Kirkpatrick model.  The cutoff location is
determined by an equilibrium gradient operator constructed in this
paper.  The distinction between equilibrium relaxation and worst-start
mixing is central: a positive spectral gap controls equilibrium
fluctuations, but a cutoff theorem must also identify when an arbitrary
initial configuration loses its remaining bias.

The Sherrington--Kirkpatrick model has random interactions of mixed
signs \cite{SK1975}.  These interactions make the required control of
nonequilibrium trajectories delicate even at high temperature.

Functional inequalities have provided quantitative high-temperature
control. The random-SK bounds in the following comparison hold with
probability tending to one over the disorder, at each fixed temperature
in the stated range.  Bauerschmidt and Bodineau~\cite{BB2019} proved a uniform
logarithmic Sobolev inequality for \(\beta<1/4\) with an unweighted
flip-gradient form.  Eldan, Koehler, and Zeitouni~\cite{EKZ2022}
obtained a heat-bath Poincar\'e inequality in that range.
Anari, Jain, Koehler, Pham, and Vuong~\cite{AJKPV2022} obtained
\(O(n\log n)\) mixing there, uniformly over external fields, through
entropic independence. Adhikari, Brennecke, Xu, and Yau~\cite{ABXY2024}
proved heat-bath gap estimates for mixed \(p\)-spin models at sufficiently
high temperature.
The localization framework of Chen and Eldan~\cite{CE2025} relates
entropy contraction to covariance estimates along tilted measures.
Anari, Koehler, and Vuong~\cite{AKV2024} extended
\(O(n\log n)\) Glauber mixing to a threshold approximately \(0.295\).
Wang~\cite{Wang2026} proved
\(O_\beta(n\log(n/\varepsilon))\) worst-start mixing for every fixed
\(\beta<1/2\), uniformly over external fields.
These mixing times count single-site update attempts.
Boban, Li, and Oveis Gharan~\cite{BLO2026} proved a positive unscaled
heat-bath gap in zero field for \(\beta<1/2+\varepsilon_0\), with a
universal \(\varepsilon_0\geq5\cdot10^{-5}\), yielding \(O_\beta(n^2)\) single-site update attempts at fixed
total-variation accuracy.
These results do not identify a cutoff location.

The companion article \cite{SKGap} proves a positive unscaled heat-bath
gap for every fixed \(\beta<1\) in the zero-field Gaussian model.
We record the precise imports in \Cref{sec:model} and prove the
additional arbitrary-law mean, path, and gradient estimates needed for
cutoff here.  They identify the leading location through a deterministic
equilibrium spectral construction.  The resulting rate is implicit.
The appendices give an independent proof of the cutoff ratio for
\(\beta<1/2\), with its own scalar certificate and auxiliary estimates
uniform over external fields and interaction scalings.  Those stronger auxiliary uniformities
are confined to the restricted-temperature method.

\subsection{The result}

For \(\Omega_n=\{-1,1\}^n\), let \(J\) be symmetric with zero diagonal
and independent \(J_{ij}\sim N(0,\beta^2/n)\) for \(i<j\). The Gibbs law is
\[
 \mu_J(x)=Z_J^{-1}\exp\!\left(\frac12x^{\mathsf T}Jx\right).
\]
This normalization puts the equilibrium high-temperature boundary at
\(\beta=1\): the normalized mean pressure \(n^{-1}\E\log Z_J\)
tends to \(\log2+\beta^2/4\) for \(\beta<1\)
\cite[Proposition~2.1]{ALR1987}. For \(\beta>1\), Guerra's
variational bound~\cite{Guerra2003} gives a limit superior strictly
smaller than \(\log2+\beta^2/4\).

At each spin a rate-one clock replaces that spin by a sample from its
conditional law under \(\mu_J\), given the other spins. Write \(S_t\)
for this transition kernel and \(P\) for one uniformly chosen single-site
refresh attempt; attempts that retain the old spin are included.
At fixed disorder put
\[
 d_n(t)=\max_x\|S_t(x,\cdot)-\mu_J\|_{\TV},\qquad
 d_n^{\rm disc}(k)=\max_x\|P^k(x,\cdot)-\mu_J\|_{\TV},
\]
where \(\|\nu-\mu\|_{\TV}=\frac12\sum_x|\nu(x)-\mu(x)|\).
We abbreviate \(\mu_J\) to \(\mu\). Section~\ref{sec:model} develops the
half-difference and Dirichlet-form conventions used in the proof.

\begin{theorem}
\label{thm:main}
For every fixed \(0\leq\beta<1\), there is a deterministic
\(\lambda(\beta)>0\) such that, with
\[
 t_n=\frac{\log n}{2\lambda(\beta)},
\]
for every \(\epsilon\in(0,1)\),
\begin{align*}
 d_n((1-\epsilon)t_n)&\longrightarrow1,
 &
 d_n((1+\epsilon)t_n)&\longrightarrow0,\\
 d_n^{\rm disc}(\lfloor(1-\epsilon)nt_n\rfloor)&\longrightarrow1,
 &
 d_n^{\rm disc}(\lceil(1+\epsilon)nt_n\rceil)&\longrightarrow0.
\end{align*}
All convergences are in probability over the disorder.
For \(\beta=0\), \(\lambda(0)=1\).
For \(0<\beta<1\), \(\lambda(\beta)\) is the equilibrium-operator
constant defined in \eqref{eq:exact-lambda}; it satisfies
\[
 0<\lambda(\beta)\leq(1-\beta)^2.
\]
\end{theorem}

The theorem answers the high-temperature cutoff question in the
zero-field Gaussian setting.  The temperature is fixed before the
dimension tends to infinity.  We do not address a temperature
approaching \(1\) with \(n\), external fields, or a cutoff window.

\subsection{The rate and the gradient mechanism}

A half-difference measures an observable's sensitivity to one spin:
\[
 d_i f(x)=\frac{f(x^{i,+})-f(x^{i,-})}{2},
 \qquad df=(d_1f,\ldots,d_nf),
\]
where \(x^{i,\pm}\) is obtained from \(x\) by fixing spin \(i\) to
\(\pm1\). We call a vector field of the form \(df\) an exact gradient.
Differentiating heat flow gives a vector semigroup \(K_t\) satisfying
\[
 dS_tf=K_tdf.
\]
Its finite-dimensional generator is derived in
Section~\ref{sec:gradient-path}. Thus decay under \(K_t\) controls
how strongly a smoothed observable can still distinguish its starting
configuration.

The scale relevant to total variation is visible in the elementary bound
\[
 \max_x g(x)-\min_x g(x)
 \le 2\sqrt n\,\sup_x\|dg(x)\|.
\]
For the upper bound, the proof makes the gradient of \(S_tf\) smaller
than \(n^{-1/2}\), uniformly for \(|f|\le1\), by combining gradient
propagation with a late-time energy bound. For the lower bound, it finds
an observable with bounded fluctuations and a mean bias that decays at
any rate just above \(\lambda\), from an initial size arbitrarily close
to \(\sqrt n\) on a logarithmic scale. These two thresholds explain
the location \(\log n/(2\lambda)\).

Two equilibrium rates can govern this decay. First, start the chain in
\(\mu\) and consider its normalized spin autocorrelation
\[
 C_n(t)=\frac1n\E_\mu[X_0^{\mathsf T}X_t].
\]
Section~\ref{sec:row-spin} proves convergence to
\(C(t)=\int e^{-ut}\rho(\dd u)\) for a deterministic positive
probability measure \(\rho\). Its lower support edge is
\(\lambda_{\rm sp}=\inf\supp\rho>0\).
This rate describes the slowest part detected by the limiting averaged
spin correlation.

The second rate is obtained by retaining more equilibrium observables.
A finite calculation starts with the spin vector and fixed independent
site seeds, and uses multiplication by \(J\), linear combinations,
and smooth componentwise functions. For example, the local-field vector
\(Jx\) and the conditional-mean vector \(\tanh(Jx)\) are such
calculations. After normalization by \(n^{-1/2}\), their limiting
equilibrium overlaps define a deterministic Hilbert space of
\emph{visible} vectors. A sequence of endpoint vectors is \emph{hidden}
when its contractions with every fixed such calculation vanish in
square mean.
Section~\ref{sec:gaussian-calculus} constructs this space and its
conditional path predictions.

For exact gradients, the natural equilibrium energy uses the conditional
spin variances \(v_i(x)=1-\tanh^2((Jx)_i)\):
\[
 \|p\|_v^2=\E_\mu\sum_i v_i p_i^2.
\]
Section~\ref{sec:exact} constructs an invariant family of equilibrium
test gradients from the finite calculations and completes it in this
weighted norm. The resulting space \(\mathcal X_v\) contains every
weighted-visible limit of exact gradients relevant to the proof.
On it, the limiting \(K_t\) is self-adjoint; write
\(\mathcal A_v=-{\rm gen}(K|_{\mathcal X_v})\). This is a preview
of the precise test-potential construction in
Definition~\ref{def:exact-spaces} and its saturation result. The rate is
\[
 \lambda=\min\{\lambda_{\rm sp},\inf\spec(\mathcal A_v)\},
 \qquad
 0<\lambda\le\lambda_{\rm sp}\le(1-\beta)^2.
\]
All these objects are constructed from equilibrium quantities at fixed
times: dimension tends to infinity before any subsequent large-time or
spectral-edge limit. If the second rate is smaller, the proof isolates a
finite-dimensional family of slower observables and uses it for the
lower bound. The formula does not identify \(\lambda\) with the upper
bound \((1-\beta)^2\).

\subsection{Main ideas}

The main difficulty is to make these equilibrium decay mechanisms
control every starting configuration for logarithmically long times.
The proof establishes fixed-time comparisons first, then makes a single
sufficiently long block uniform before iterating it.

\paragraph{Bringing arbitrary starts into the comparison regime.}
We strengthen the stable-field residual argument underlying the companion
gap theorem to control the mean under an arbitrary change of law. The
change in mean is bounded by \(C(1+\sqrt D)\), where \(D\) is the
Dirichlet energy of the square-root density. Transport through a two-stage Gaussian observation
then gives entropy control and the fixed-time heat-kernel estimate
\[
 \frac1n\log\left\|\frac{S_d(x,\cdot)}{\mu}\right\|_\infty
 \le C e^{-c d^{2/3}}.
\]
Here \(d\) is fixed before \(n\to\infty\). Choosing a sufficiently
long fixed burn-in makes this density cost small enough for exponentially
reliable equilibrium comparisons to survive the changed law. The
observation argument also supplies the posterior and cavity gaps used
later; the zero-field gap alone would not supply them.

\paragraph{Averaging the hidden part from both endpoints.}
The path expansion of \(K_t\) involves products of \(J\) and diagonal
labels depending on observed spin states. We compare each label with
its conditional prediction from the initial or terminal state. A
conditional-covariance argument shows that terms containing both an
initially centered label and a later terminally centered label vanish
when tested against a hidden initial vector and an exact terminal
gradient. This leaves strings predicted from their two endpoints, with
at most one actual path label between them.

Two estimates make this averaging useful. A projected row likelihood
has subexponential moments in the interval length, while the spin
spectral measure supplies summable positive path weights at every rate
below \(\lambda_{\rm sp}\). Their combination proves hidden-gradient
decay at that edge. To justify the covariance cancellation for multitime
labels, we propagate mixed differences through conditional evolution and
copied observation states. The resulting directed networks retain the
derivative continuing to the terminal gradient. This is the role of the
derivative machinery in Sections~\ref{sec:gradient-path} and
\ref{sec:gaussian-calculus}. Its operator estimates also retain rare
coordinate types, whose small spatial frequency need not make their
operator action small.

\paragraph{Recovering the visible exact part.}
Projecting a nonstationary gradient onto the visible calculations gives
random scalar coefficients. Transport acts on the deterministic vector
directions by \(K_t\) and on these coefficients by conditional
expectation. The projection does not immediately preserve the identities
that characterize a gradient. Reverse entropy estimates and strong
control of weighted curl restore those identities; a second observation
argument then bounds the distance to actual gradients. This places the
visible projection in \(\mathcal X_v\) without differentiating its
possibly irregular coefficients. The hidden bound further shows that
any visible spectrum below \(\lambda_{\rm sp}\) consists of isolated
modes of finite multiplicity.

\paragraph{From a fixed block to the two cutoff bounds.}
The hidden and visible estimates give contraction on one long fixed
block. Strict contraction slack and normalization make this bound uniform
over all starts and tests, with an inverse-polynomial additive error;
iteration therefore reaches logarithmic times without assuming a
quantitative rate for a fixed-time limit. A late-time martingale-energy
estimate completes the upper bound. For the lower bound, spin correlation
provides the distinguishing observable when
\(\lambda=\lambda_{\rm sp}\). If a slower visible mode determines
\(\lambda\), an additional block estimate aligns both the mean and
fluctuations of its random coefficient. A small perturbation of the
initial law then creates a detectable slow bias, and convexity selects
a deterministic starting configuration with at least the same
separation from equilibrium.

\subsection{Organization and dependencies}

Section~\ref{sec:model} fixes the normalizations and companion inputs,
and states the finite-chain concentration and clock estimates.
Section~\ref{sec:static} proves the arbitrary-law and observation bounds.
Sections~\ref{sec:gradient-path} and~\ref{sec:gaussian-calculus} develop
the gradient evolution, conditional derivatives, and Gaussian path
calculus. Section~\ref{sec:row-spin} treats the projected row likelihood
and the spin spectral measure; Section~\ref{sec:hidden} uses them to
prove hidden propagation. Section~\ref{sec:exact} constructs the visible
exact operator and isolates its slower spectrum. Section~\ref{sec:cutoff}
assembles the uniform blocks and proves cutoff for both clocks.
Appendices~\ref{sec:rcut-introduction}--\ref{rcut:sec:scalar-certificate}
give the independent restricted-temperature argument and its scalar
certificate. Appendix~\ref{app:clock-comparison} proves the common
clock lemma. Figure~\ref{fig:proof-dependencies} records the principal
dependencies.

\begin{figure}[htbp]
\centering
\begin{tikzpicture}[
  box/.style={draw=black!60,rounded corners=2pt,fill=black!3,
    text width=5.15cm,align=center,inner sep=7pt,font=\small},
  arrow/.style={-{Stealth[length=2.2mm]},draw=black!65,line width=.6pt},
  node distance=11mm and 13mm]
 \node[box] (static) {Observation and entropy (\S\ref{sec:static})\\
   Density and posterior control};
 \node[box,right=of static] (paths) {Finite-time Gaussian calculus\\(\S\S\ref{sec:gradient-path}--\ref{sec:gaussian-calculus})\\
   Conditional predictors\\and matrix words};
 \node[box,below=of static] (row) {Projected row moments (\S\ref{sec:row-spin})\\
   Spin spectral edge \(\lambda_{\rm sp}\)};
 \node[box,below=of paths] (exact) {Weighted gradient exactness (\S\ref{sec:exact})\\
   The visible sector \(\mathcal X_v\)};
 \node[box,below=of row] (hidden) {Two-ended path averaging (\S\ref{sec:hidden})\\
   Hidden decay at \(\lambda_{\rm sp}\)};
 \node[box,below=of exact] (spectral) {Spectral isolation (\S\ref{sec:exact})\\
   Slower visible eigenmodes};
 \coordinate (middle) at ($(hidden.south)!.5!(spectral.south)$);
 \node[box,below=14mm of middle,
   text width=8cm] (cutoff)
   {Uniform fixed-block estimates (\S\ref{sec:cutoff})\\
    Upper cutoff and the two spectral lower bounds};
 \draw[arrow] (static) -- (row);
 \draw[arrow] (paths) -- (row);
 \draw[arrow] (paths) -- (exact);
 \draw[arrow] (static) -- (exact);
 \draw[arrow] (row) -- (hidden);
 \draw[arrow] (exact) -- (spectral);
 \draw[arrow] (hidden) -- (spectral);
 \draw[arrow] (hidden) -- (cutoff);
 \draw[arrow] (spectral) -- (cutoff);
\end{tikzpicture}

\caption{The principal estimates and their use in the cutoff proof.
All limiting comparisons precede the iteration over logarithmically
many fixed-time blocks. The stationary gap is an input; the
reweighting, averaging, and exactness arguments are proved here.}
\label{fig:proof-dependencies}
\end{figure}
